% GENERATED FILE. Edit canonical lessons, then run npm run generate:books.
% canonical-source-hash: fnv1a64:788a3d766290bb1f
% canonical-lessons: TE-C411-janaganana, TE-C411-sarve, TE-C411-ganankalu, TE-C411-prajabhiprayam, TE-C411-satam

\chapter{Counting People}
\label{ch:te-counting-people}
\hlchaptermodality{411}

\begin{chapteropening}
\textbf{By the end of this chapter:} \emph{I can talk about censuses, surveys, figures and percentages.}

Complete the last lesson of chapter 411: I can talk about censuses, surveys, figures and percentages.
\end{chapteropening}

\section[\texorpdfstring{\te{జనగణన} (janagaṇana)}{జనగణన (janagaṇana)}]{\te{జనగణన} (janagaṇana) — a census}

\label{lesson:TE-C411-janaganana}



\begin{quote}
Before the new one: say the Telugu for an activist, then the Telugu for a donor.
\end{quote}

\subsection*{You'll want to know: \te{జనగణన}}
\textbf{\te{జనగణన}} — \emph{janagaṇana} — \textquotedblleft{}a census\textquotedblright{}.

\begin{quote}
\te{పది} \te{సంవత్సరాలకు} \te{ఒకసారి} \te{జనగణన} \te{జరుగుతుంది}.
\end{quote}

\begin{quote}
\emph{padi saṁvatsarālaku okasāri janagaṇana jarugutundi.}
\end{quote}

\begin{quote}
A census takes place once every ten years.
\end{quote}

\te{జన}, people, + \te{గణన}, counting.

\begin{cousinweb}
\textbf{\te{జనగణన}} (\emph{janagaṇana}) is built from Sanskrit \emph{jana}, \textquotedblleft{}a person, people\textquotedblright{} (CDIAL 5098), and \emph{gaṇana}, \textquotedblleft{}counting\textquotedblright{} (CDIAL 3992).
\end{cousinweb}

\begin{grammarlens}[title={What you've built}]
One more word for giving reasons and opinions.
\end{grammarlens}

\subsection*{Your turn}
\begin{itemize}
\raggedright
  \item \emph{Say it:} \emph{janagaṇana}
  \item \emph{Say it:} \emph{janagaṇana}, once more
  \item \emph{Say it:} \emph{padi saṁvatsarālaku okasāri janagaṇana jarugutundi.} — A census takes place once every ten years.
  \item \emph{Recall:} say the Telugu for an activist, then the Telugu for a donor, then say \emph{janagaṇana} again
\end{itemize}

\begin{tcolorbox}[breakable,colback=teal!4,colframe=teal!35!black,title={Before you move on}]
What does \te{జనగణన} mean? (\textquotedblleft{}A census\textquotedblright{}.) Say it once more.
\end{tcolorbox}
\section[\texorpdfstring{\te{సర్వే} (sarvē)}{సర్వే (sarvē)}]{\te{సర్వే} (sarvē) — a survey}

\label{lesson:TE-C411-sarve}



\begin{quote}
Before the new one: say the Telugu for a donor, then the Telugu for a census.
\end{quote}

\subsection*{You'll want to know: \te{సర్వే}}
\textbf{\te{సర్వే}} — \emph{sarvē} — \textquotedblleft{}a survey\textquotedblright{}.

\begin{quote}
\te{ఆ} \te{సర్వే} \te{ప్రకారం} \te{జనం} \te{ధరల} \te{గురించి} \te{బాధపడుతున్నారు}.
\end{quote}

\begin{quote}
\emph{ā sarvē prakāraṁ janaṁ dharala guriñci bādhapaḍutunnāru.}
\end{quote}

\begin{quote}
According to that survey, people are worried about prices.
\end{quote}

\te{సర్వే} \te{చేయు}, to carry out a survey.

\begin{cousinweb}
\textbf{\te{సర్వే}} (\emph{sarvē}) is a loan from English \emph{survey}.
\end{cousinweb}

\begin{grammarlens}[title={What you've built}]
One more word for giving reasons and opinions.
\end{grammarlens}

\subsection*{Your turn}
\begin{itemize}
\raggedright
  \item \emph{Say it:} \emph{sarvē}
  \item \emph{Say it:} \emph{sarvē}, once more
  \item \emph{Say it:} \emph{ā sarvē prakāraṁ janaṁ dharala guriñci bādhapaḍutunnāru.} — According to that survey, people are worried about prices.
  \item \emph{Recall:} say the Telugu for a donor, then the Telugu for a census, then say \emph{sarvē} again
\end{itemize}

\begin{tcolorbox}[breakable,colback=teal!4,colframe=teal!35!black,title={Before you move on}]
What does \te{సర్వే} mean? (\textquotedblleft{}A survey\textquotedblright{}.) Say it once more.
\end{tcolorbox}
\section[\texorpdfstring{\te{గణాంకాలు} (gaṇāṅkālu)}{గణాంకాలు (gaṇāṅkālu)}]{\te{గణాంకాలు} (gaṇāṅkālu) — statistics; figures}

\label{lesson:TE-C411-ganankalu}



\begin{quote}
Before the new one: say the Telugu for a census, then the Telugu for a survey.
\end{quote}

\subsection*{You'll want to know: \te{గణాంకాలు}}
\textbf{\te{గణాంకాలు}} — \emph{gaṇāṅkālu} — \textquotedblleft{}statistics; figures\textquotedblright{}.

\begin{quote}
\te{ఆ} \te{గణాంకాలు} \te{సరైనవి} \te{కావు}.
\end{quote}

\begin{quote}
\emph{ā gaṇāṅkālu sarainavi kāvu.}
\end{quote}

\begin{quote}
Those figures are not correct.
\end{quote}

A plural noun: \te{గణ}, a number, + \te{అంకాలు}, marks, numerals.

\begin{cousinweb}
\textbf{\te{గణాంకాలు}} (\emph{gaṇāṅkālu}) is built from Sanskrit \emph{gaṇa}, \textquotedblleft{}a troop, a number\textquotedblright{} (CDIAL 3988), and \emph{aṅka}, \textquotedblleft{}a mark, a sign\textquotedblright{} (CDIAL 100); Telugu uses the plural for statistics.
\end{cousinweb}

\begin{grammarlens}[title={What you've built}]
One more word for giving reasons and opinions.
\end{grammarlens}

\subsection*{Your turn}
\begin{itemize}
\raggedright
  \item \emph{Say it:} \emph{gaṇāṅkālu}
  \item \emph{Say it:} \emph{gaṇāṅkālu}, once more
  \item \emph{Say it:} \emph{ā gaṇāṅkālu sarainavi kāvu.} — Those figures are not correct.
  \item \emph{Recall:} say the Telugu for a census, then the Telugu for a survey, then say \emph{gaṇāṅkālu} again
\end{itemize}

\begin{tcolorbox}[breakable,colback=teal!4,colframe=teal!35!black,title={Before you move on}]
What does \te{గణాంకాలు} mean? (\textquotedblleft{}Statistics; figures\textquotedblright{}.) Say it once more.
\end{tcolorbox}
\section[\texorpdfstring{\te{ప్రజాభిప్రాయం} (prajābhiprāyaṁ)}{ప్రజాభిప్రాయం (prajābhiprāyaṁ)}]{\te{ప్రజాభిప్రాయం} (prajābhiprāyaṁ) — public opinion}

\label{lesson:TE-C411-prajabhiprayam}



\begin{quote}
Before the new one: say the Telugu for a survey, then the Telugu for statistics.
\end{quote}

\subsection*{You'll want to know: \te{ప్రజాభిప్రాయం}}
\textbf{\te{ప్రజాభిప్రాయం}} — \emph{prajābhiprāyaṁ} — \textquotedblleft{}public opinion\textquotedblright{}.

\begin{quote}
\te{ప్రజాభిప్రాయం} \te{ఆ} \te{ప్రతిపాదనకు} \te{వ్యతిరేకంగా} \te{ఉంది}.
\end{quote}

\begin{quote}
\emph{prajābhiprāyaṁ ā pratipādanaku vyatirēkaṅgā undi.}
\end{quote}

\begin{quote}
Public opinion is against that proposal.
\end{quote}

\te{ప్రజా}, of the people, + \te{అభిప్రాయం}, an opinion, which you know.

\begin{cousinweb}
\textbf{\te{ప్రజాభిప్రాయం}} (\emph{prajābhiprāyaṁ}) is built from Sanskrit \emph{prajā}, \textquotedblleft{}offspring, people\textquotedblright{}, and \emph{abhiprāya}, \textquotedblleft{}aim, intention\textquotedblright{} (Monier-Williams), the word Telugu has as \te{అభిప్రాయం}.
\end{cousinweb}

\begin{grammarlens}[title={What you've built}]
One more word for giving reasons and opinions.
\end{grammarlens}

\subsection*{Your turn}
\begin{itemize}
\raggedright
  \item \emph{Say it:} \emph{prajābhiprāyaṁ}
  \item \emph{Say it:} \emph{prajābhiprāyaṁ}, once more
  \item \emph{Say it:} \emph{prajābhiprāyaṁ ā pratipādanaku vyatirēkaṅgā undi.} — Public opinion is against that proposal.
  \item \emph{Recall:} say the Telugu for a survey, then the Telugu for statistics, then say \emph{prajābhiprāyaṁ} again
\end{itemize}

\begin{tcolorbox}[breakable,colback=teal!4,colframe=teal!35!black,title={Before you move on}]
What does \te{ప్రజాభిప్రాయం} mean? (\textquotedblleft{}Public opinion\textquotedblright{}.) Say it once more.
\end{tcolorbox}
\section[\texorpdfstring{\te{శాతం} (śātaṁ)}{శాతం (śātaṁ)}]{\te{శాతం} (śātaṁ) — a percentage; per cent}

\label{lesson:TE-C411-satam}



\begin{quote}
Before the new one: say the Telugu for statistics, then the Telugu for public opinion.
\end{quote}

\subsection*{You'll want to know: \te{శాతం}}
\textbf{\te{శాతం}} — \emph{śātaṁ} — \textquotedblleft{}a percentage; per cent\textquotedblright{}.

\begin{quote}
\te{మా} \te{తరగతిలో} \te{తొంభై} \te{శాతం} \te{మంది} \te{పాసయ్యారు}.
\end{quote}

\begin{quote}
\emph{mā taragatilō tombhai śātaṁ mandi pāsayyāru.}
\end{quote}

\begin{quote}
Ninety per cent of our class passed.
\end{quote}

With numbers spelt out: \te{యాభై} \te{శాతం}, fifty per cent.

\begin{cousinweb}
\textbf{\te{శాతం}} (\emph{śātaṁ}) goes back to Sanskrit \emph{śata}, \textquotedblleft{}a hundred\textquotedblright{} (CDIAL 12278): so many out of a hundred.
\end{cousinweb}

\begin{grammarlens}[title={What you've built}]
One more word for giving reasons and opinions.
\end{grammarlens}

\subsection*{Your turn}
\begin{itemize}
\raggedright
  \item \emph{Say it:} \emph{śātaṁ}
  \item \emph{Say it:} \emph{śātaṁ}, once more
  \item \emph{Say it:} \emph{mā taragatilō tombhai śātaṁ mandi pāsayyāru.} — Ninety per cent of our class passed.
  \item \emph{Recall:} say the Telugu for statistics, then the Telugu for public opinion, then say \emph{śātaṁ} again
\end{itemize}

\begin{tcolorbox}[breakable,colback=teal!4,colframe=teal!35!black,title={Before you move on}]
What does \te{శాతం} mean? (\textquotedblleft{}A percentage; per cent\textquotedblright{}.) Say it once more.
\end{tcolorbox}
